\section{Dataset and Split Protocol}
\label{sec:data}

% Budget: about 0.6 column-pages. Rubric: split protocol is worth 5 implementation points
% ("correct unit identified, both splits run, test set used once").

%-------------------------------------------------------------------------
\subsection{Dataset}
% Source: distracted_driver.ipynb Section 1, fig01_class_and_driver_balance, fig02_class_examples.
% 22,424 labeled frames, 26 drivers, 10 classes, one fixed camera, 640x480.
% Classes are close to balanced but images per driver vary.
% c1/c3 and c2/c4 differ only by hand, which is why plain horizontal flips corrupt labels.
% Data use (R12): one sentence that we follow the Kaggle competition rules, do not
% redistribute images, and checked the rules before showing frames in this paper.
\TODO{Dataset statistics, class list, and data use statement. Optional small figure of one frame per class.}

%-------------------------------------------------------------------------
\subsection{Split Protocol}
% Source: distracted_driver.ipynb Sections 1 and 6.1 (table_folds.csv).
% The paragraphs below are the earlier draft. Update them to the final protocol:
%   unit = driver. 4 drivers frozen by seed (SEED = 42) as the test set, never used for
%   training, curves or choices. The other 22 development drivers are split with GroupKFold
%   into 3 driver-disjoint folds. Each final model is trained on all 22 and scored once.
% If the interim checkpoint declared a different split (the 70/30 split below), say so and
% say why it changed. A silent change looks like test-set shopping.
% State explicitly that no model or setting was chosen from test-set numbers.

All data was split based on the provided 26 drivers, using 70\% of the set for training, and the final 30\% for testing validation.
The split was based on previous industry reports\cite{maDigitalTwin2024} on the best split for their training model and became the basis for ours.
A random image level split allows any duplicate frames of the same driver taken from the original videos in both the training and testing data for confirmation.
Using \verb+GroupShuffleSplit+ with \verb+subject+ as the group, 200 candidate seeds with the split seed DEFINE created a held-out fraction near 30\% where all ten classes appear and chosen for performance.
The final split uses NUMBER drivers for training (15,703 images) and NUMBER drivers for testing (6,721 images) with confirmation that the sets are disjoint.

%-------------------------------------------------------------------------
\subsection{Metrics}
% Source: table_results_per_fold.csv columns. Defend each choice (R6, results rubric 4 pts).
%   Macro-F1 as the primary score (balanced classes, but per-class errors differ).
%   Accuracy for comparison with prior work. Multi-class log-loss, the Kaggle metric,
%   which also penalizes over-confidence (a calibration signal). ROC-AUC on the test set.
%   Per-class F1 and confusion matrices to show where errors sit.
%   Driver-identity probe, a linear classifier that names the held-out driver from frozen
%   features, compared against chance (1 / number of held-out drivers). It measures the
%   mechanism, not the task.
\TODO{Metrics and why each one was chosen.}
